\documentclass[10pt]{ltjsarticle}
\usepackage{tem}
\usepackage{hyperref}
\usetikzlibrary{positioning,arrows.meta}
\pagestyle{fancy}
\lhead{Infinite Generation Two}
\rhead{Pi extension specification}
\setlength{\headheight}{16pt}
\setlength{\headsep}{30pt}
\begin{document}
\title{Infinite Generation Two\\\large{--- Native Pi Extension ---}}
\author{Fides Thérèse}
\date{2026-10-22}
\maketitle
\thispagestyle{plain}

\section{Goal and acceptance criteria}
The repository becomes a native Pi extension distributed from
\url{https://github.com/FidesTherese/dsh-infinite-gen-2}. The extension MUST
preserve the exact bytes of \texttt{prompts/infinite-gen-2.md}, expose the
existing \texttt{infinite\_gen2\_profile} tool, append the bundled instructions
without replacing Pi's base instructions, and display a session-local status
indicator. The README MUST be English, contain accurate Pi installation
instructions, and omit advertisements, donation requests, and promotional
community links. The MIT notice and upstream attribution remain intact.

Acceptance requires automated tests, strict TypeScript checking, a real Pi
loader/event-runner smoke test, and a package-content check. The requested
handoff includes a reviewed commit pushed to the fork's existing
\texttt{master} branch without rewriting history.

\section{Original repository analysis}
The inspected baseline is commit \texttt{3a701c3}. This is a small ES-module
DeepSeek Harness plugin, not a standalone model client. Its execution path is
\texttt{apply(ctx)} in \texttt{index.js}: read the bundled Markdown, register an
ordered system-prompt section, register a profile tool, and optionally register
an in-memory response-status projection. \texttt{client.js} consumes that
projection through a host browser module loader and React, then inserts a
badge into the conversation input dock. The Cordis patch declares the plugin
bundle. Four shell/PowerShell scripts copy files and alter Harness profile
configuration; the Windows installer also registers a desktop URI scheme.
There are no application dependencies or automated tests declared at baseline.

The host APIs, browser globals, and profile installers do not exist in Pi.
The repository's reusable assets are the prompt, profile metadata, and the
bilingual text-marker heuristic. The README's numeric strength comparisons
are inherited calibration labels, not demonstrated benchmarks. The detector
can report quoted phrases as refusals and cannot measure model compliance.

Recommended reading order is package metadata, the original entry point,
the prompt asset, the projection scorer, and finally the browser/installer
adapters. Git preserves these removed adapters at the baseline commit.

\section{Chosen design and module boundaries}
A direct TypeScript port is preferred to emulating Harness APIs or maintaining
two host adapters: Pi already supplies prompt sections, tools, events, and
status UI. \texttt{index.ts} owns Pi registration and ephemeral session state.
\texttt{src/profile.ts} owns asset integrity and the read-only profile-tool
contract. \texttt{src/armor.ts} owns pure text extraction, marker matching,
and English status formatting. These files form one cohesive extension
module; no service, database, browser, or provider integration is introduced.

The following nodes are new or replaced in this port. The prompt node is
unchanged. Dependencies point from the adapter to its helpers, never back.
\begin{center}
\definecolor{TransBlue}{HTML}{5BCEFA}
\definecolor{TransPink}{HTML}{F5A9B8}
\begin{tikzpicture}[font=\small, node distance=9mm and 12mm,
  box/.style={draw,rounded corners,text width=.24\linewidth,align=center,
  minimum height=11mm,fill opacity=.55,text opacity=1}, >=Stealth]
  \node[box,fill=TransBlue] (pi) {Pi host\\events, tools, UI};
  \node[box,fill=TransPink,right=of pi] (entry) {Native adapter\\index.ts};
  \node[box,fill=TransBlue,below=of pi] (armor) {Pure status helper\\src/armor.ts};
  \node[box,fill=TransPink,below=of entry] (profile) {Profile and integrity\\src/profile.ts};
  \node[box,fill=white,below=of profile] (prompt) {Unchanged prompt\\Markdown asset};
  \draw[->] (pi) -- (entry);
  \draw[->] (entry) -- (armor);
  \draw[->] (entry) -- (profile);
  \draw[->] (profile) -- (prompt);
\end{tikzpicture}
\end{center}

\section{Interfaces, state, and critical flow}
Pi discovers \texttt{index.ts} through the explicit package manifest. On load,
the adapter verifies the prompt's SHA-256 and registers the tool, command,
flag, and event handlers. Before each run, the adapter sets only its own
\texttt{infinite\_gen2} prompt section. Repeating the event replaces the same
section rather than duplicating text; disabling removes only that section.
It does not replace Pi's tools, project context, provider, or credentials.

The parameterless profile tool returns the complete unchanged prompt, its
fingerprint, upstream metadata, prompt version 0.2.3, and extension version
0.3.0 as model-facing text, details, and structured JSON. The profile remains
available while prompt injection is disabled.

The extension starts enabled. \texttt{--infinite-gen2-off} disables it at
session startup. \texttt{/infinite-gen2 on|off|status} changes or reports the
current session state; an empty argument means status. Commands never start
an additional model turn. State resets on session startup/reload/replacement;
a tree move clears the verdict. It is not durable configuration.

\texttt{agent\_start} clears old results and marks the run busy.
\texttt{message\_end} inspects only assistant text, excluding thinking and
tool blocks. Errors and cancellations, including automatic compaction failures,
receive distinct outcomes rather than passing. Toggle commands preserve the
busy state of an active run. \texttt{agent\_settled}, not the low-level \texttt{agent\_end}, ends
the busy indicator after retries and queued work. Shutdown clears the status.
The UI receives status only when available; headless model/tool behavior is
identical and produces no unsolicited stdout.

\section{Constraints, risks, and non-goals}
The target is Node.js 22.19 or newer and Pi 0.99.1 or a compatible newer
release. TypeScript sources load directly through Pi; no compilation is
required for installation. Pi and TypeBox are host-provided peer dependencies.
Development versions are pinned in the lockfile.

A missing or changed prompt makes extension loading fail visibly; the port
never substitutes another prompt. The frozen SHA-256 fingerprint is:
\par\noindent\texttt{02679eb726507d5a6bfe0a327c8469965440be7a7c5104cab57b8d1e529ec466}
\par
Other extensions that force-replace the entire system prompt can override
structured sections. Status scoring is an inherited text heuristic, not a
security or effectiveness guarantee. The extension changes model instructions
but does not alter provider safeguards, OS permissions, or tool approvals.

DeepSeek Harness compatibility, browser animations, URI-scheme installers,
persistent toggles, automatic retries based on verdicts, telemetry, and prompt
translation are out of scope. The obsolete Harness code is removed rather
than advertised as supported. Implementation detail and the dependency-safe
roadmap belong in \texttt{docs/DTD.tex}; module context belongs in
\texttt{docs/modules/pi-extension.tex}.

\section{Verification record}
Verification on 2026-10-22: \texttt{npm test} exercises 31 focused SDK/unit
cases; \texttt{npm run check} performs strict TypeScript checking;
\texttt{npm run smoke} exercises three offline CLI/package-install cases;
\texttt{npm pack --dry-run} confirms seven required distribution files.
The package loads through the Pi 0.99.1 CLI and was also checked with the
current Pi 1.4.0 CLI. Tests use isolated temporary Pi state and make no model
requests. Live DeepSeek behavior is not an acceptance claim.
\end{document}
